% =====================================================================
% Appendix: proofs for Section 3 (Search makes soundness a worst-case property)
% Sources: research/theory/T1-soundness-under-search.md (Cor 2.2, Prop 2.3 and its repair),
%          research/lit/L3-learning-to-reason-and-limit-inquiry.md (Lemma A, as repaired),
%          research/theory/T2-coherence-as-negative-data.md (Thm 2.4, as repaired),
%          research/verification/{T1,T2,L3}-*.md.
% =====================================================================

\section{Proofs for Section~\ref{sec:search}}\label{app:search}

\Cref{lem:search:stepwise}, \cref{thm:search:tonk}, \cref{prop:search:post} and \cref{prop:search:mdl} are proved in full in the main text. This appendix proves the remaining results of \cref{sec:search} and verifies the claims of \cref{ex:search:atoms}. Notation is as in \cref{sec:search}. Throughout, $\models$ is classical propositional consequence, and a step $\Pi\step\varphi$ over formulas is \emph{sound} if $\Pi\models\varphi$.

% ---------------------------------------------------------------------
\subsection{PAC learners can be maximally unsound}\label{app:search:pac}

\begin{proof}[Proof of \cref{cor:search:pac}]
For $N\ge0$ let $\mathcal S_{\ge N}:=\{s\in\mathcal S:|s|\ge N\}$.

\emph{Step 1: $T_j\subseteq\mathcal S_{\ge j}$.} Every step of~$T_j$ has the premise $A\vee\top^{(j)}$, and $\top^{(j)}$ contains $j$ occurrences of~$\top$. So every step of~$T_j$ has size at least~$j$.

\emph{Step 2: the tail bound.} The sets $\mathcal S_{\ge N}$ decrease in~$N$ and have empty intersection, so $Q(\mathcal S_{\ge N})\downarrow0$ by continuity from above. Fix $\varepsilon\in(0,1)$ and let $N_\varepsilon:=\min\{N\ge0:Q(\mathcal S_{\ge N})\le\varepsilon\}$.
\begin{itemize}
\item If $\hat\jmath\ge N_\varepsilon$, then $Q(T_{\hat\jmath})\le Q(\mathcal S_{\ge\hat\jmath})\le Q(\mathcal S_{\ge N_\varepsilon})\le\varepsilon$, by Step~1 and monotonicity.
\item If $\hat\jmath<N_\varepsilon$, then $N_\varepsilon\ge2$, because $\hat\jmath\ge1$. By minimality of~$N_\varepsilon$, $Q(\mathcal S_{\ge N_\varepsilon-1})>\varepsilon$. Moreover $\hat\jmath<N_\varepsilon$ means $|X_i|\le N_\varepsilon-2$ for every~$i$, i.e.\ no sample lies in $\mathcal S_{\ge N_\varepsilon-1}$. Since the $X_i$ are i.i.d.\ from~$Q$, this event has probability $(1-Q(\mathcal S_{\ge N_\varepsilon-1}))^m<(1-\varepsilon)^m$.
\end{itemize}
Hence $\Pr[Q(T_{\hat\jmath})>\varepsilon]\le\Pr[\hat\jmath<N_\varepsilon]\le(1-\varepsilon)^m$.

\emph{Step 3: the error decomposition.} Write $h=\mathcal L(X_{1..m})$ and $Y=\inst(\top\mathrm{I},\wedge\mathrm{I},\vee\mathrm{I}_2)\subseteq\Rstar$. An element of $(h\cup T_{\hat\jmath}\cup Y)\setminus\Rstar$ lies in $h\setminus\Rstar$ or in~$T_{\hat\jmath}$, because $Y\subseteq\Rstar$. An element of $\Rstar\setminus(h\cup T_{\hat\jmath}\cup Y)$ lies in $\Rstar\setminus h$. So $(h\cup T_{\hat\jmath}\cup Y)\,\triangle\,\Rstar\subseteq(h\,\triangle\,\Rstar)\cup T_{\hat\jmath}$, and $\mathrm{err}_Q(\mathcal L^+)\le\mathrm{err}_Q(\mathcal L)+Q(T_{\hat\jmath})$ pointwise in the sample. This proves~(a).

\emph{Step 4: PAC.} Let $m\ge\max\{m_{\mathcal L}(\varepsilon/2,\delta/2),\lceil(2/\varepsilon)\ln(2/\delta)\rceil\}$. Then $\Pr[\mathrm{err}_Q(\mathcal L)>\varepsilon/2]\le\delta/2$, since as usual the guarantee of~$\mathcal L$ holds for every sample size $m\ge m_{\mathcal L}(\varepsilon,\delta)$. By Step~2 applied with $\varepsilon/2$, $\Pr[Q(T_{\hat\jmath})>\varepsilon/2]\le(1-\varepsilon/2)^m\le e^{-m\varepsilon/2}\le\delta/2$. By Step~3 and a union bound, $\Pr[\mathrm{err}_Q(\mathcal L^+)>\varepsilon]\le\delta$. Both events are computed from the same sample, so no extra samples beyond the maximum are needed. The outputs of~$\mathcal L^+$ need not lie in the hypothesis class of~$\mathcal L$, so the guarantee is improper. This proves~(b).

\emph{Step 5: trivialization.} Every output of $\mathcal L^+$ contains $Y\cup T_{\hat\jmath}$. The derivation in the proof of \cref{thm:search:tonk}(i), with $j=\hat\jmath$, uses only steps of~$Y$ (namely $\top$I, $\wedge$I and $\vee$I$_2$) and one step of~$T_{\hat\jmath}$. So $\Cl_{\mathcal L^+(X_{1..m})}(\emptyset)=\Fm$ for every sample. This proves~(c).
\end{proof}

\begin{remark}\label{rem:app:search:pac}
Step~5 is where the three added schemas are needed. If $\mathcal L\equiv\emptyset$ and they are omitted, the output is $T_{\hat\jmath}$ alone. Every step of~$T_{\hat\jmath}$ has a premise, so $\emptyset$ is closed under it and $\Cl_{T_{\hat\jmath}}(\emptyset)=\emptyset$.
\end{remark}

% ---------------------------------------------------------------------
\subsection{Schemas that mention atoms, and purification}\label{app:search:purify}

\paragraph{Verification of \cref{ex:search:atoms}.}
Let $\tau=(x\vee p_0\step x)$ and $A=\Rstar_{\CPC}\cup\inst(\tau)$.
\begin{enumerate}[label=(\alph*)]
\item \emph{$\tau$ is unsound.} Its instance $p_1\vee p_0\step p_1$ fails under $v(p_0)=1$, $v(p_1)=0$.
\item \emph{$A$ is not trivial.} Let $X:=\Cn_{\CPC}(\neg p_0)$. Then $X\supseteq\emptyset$. $X$ is closed under $\Rstar_{\CPC}$, because $\Cl_{\Rstar_{\CPC}}(X)=\Cn_{\CPC}(X)=X$. $X$ is also closed under $\inst(\tau)$: if $B\vee p_0\in X$, every valuation with $v(p_0)=0$ satisfies $B\vee p_0$ and hence~$B$, so $\neg p_0\models B$ and $B\in X$. By leastness, $\Cl_A(\emptyset)\subseteq X$. Since $p_0\notin X$ (take $v(p_0)=0$), $\Cl_A(\emptyset)\ne\Fm$.
\item \emph{$A$ is unsound.} The tautology $\neg p_0\vee p_0$ lies in $\Cl_{\Rstar_{\CPC}}(\emptyset)$, and the instance of~$\tau$ with $x:=\neg p_0$ yields $\neg p_0$, which is not a tautology.
\item \emph{$\inst(\tau)$ is not substitution-closed.} The substitution $p_0\mapsto p_1$ maps the instance $p_2\vee p_0\step p_2$ to $p_2\vee p_1\step p_2$, which is not an instance of~$\tau$.
\item \emph{A ground error does not fire.} For the single step $p_0\step p_1$, the set $\Cn_{\CPC}(\emptyset)$ is closed under $\Rstar_{\CPC}\cup\{p_0\step p_1\}$, because $p_0\notin\Cn_{\CPC}(\emptyset)$. So the closure from~$\emptyset$ is just the set of tautologies.
\end{enumerate}

\begin{proof}[Proof of \cref{lem:search:purify}]
Let $\tau$ have metavariables $\bar x=x_1,\dots,x_k$, and let $p_{i_1},\dots,p_{i_l}$ be the distinct atoms occurring in~$\tau$. Let $y_1,\dots,y_l$ be fresh metavariables and $\tau^\circ:=\tau[p_{i_t}:=y_t]_{t\le l}$. Write $\tau[\bar x:=\bar A]$ for the instance of~$\tau$ obtained by substituting the formulas~$\bar A$ for~$\bar x$, and similarly for~$\tau^\circ$.

($\Leftarrow$) Every instance of~$\tau$ is an instance of~$\tau^\circ$: $\tau[\bar x:=\bar A]=\tau^\circ[\bar x:=\bar A,\ \bar y:=(p_{i_1},\dots,p_{i_l})]$. So if $\tau^\circ$ is sound, so is~$\tau$.

($\Rightarrow$) Assume $\tau$ is sound, and let $\iota=\tau^\circ[\bar x:=\bar A,\ \bar y:=\bar B]$ be an instance of~$\tau^\circ$. Then $\iota$ is $\tau$ with each occurrence of~$p_{i_t}$ \emph{in $\tau$ itself} replaced by~$B_t$, and each~$x_s$ replaced by~$A_s$. Occurrences of $p_{i_t}$ inside the~$A_s$ are left alone. There are infinitely many atoms, so we can choose distinct atoms $r_1,\dots,r_l$ that occur nowhere in $\bar A$, $\bar B$ or~$\tau$. Let $\rho$ be the renaming $p_{i_t}\mapsto r_t$ ($t\le l$), identity on all other atoms, and put $A'_s:=\rho(A_s)$. Then $\iota':=\tau[\bar x:=\bar A']$ is an instance of~$\tau$, hence sound.

Let $\sigma$ be the simultaneous uniform substitution with $\sigma(p_{i_t})=B_t$, $\sigma(r_t)=p_{i_t}$ ($t\le l$), and identity on all other atoms. We claim $\sigma(\iota')=\iota$. The formulas of~$\iota'$ are built from the skeleton of~$\tau$, whose atoms are among the~$p_{i_t}$ and include no~$r_t$, with the fillers~$A'_s$ plugged in for the metavariables.
\begin{itemize}
\item On the skeleton, $\sigma$ replaces each $p_{i_t}$ by~$B_t$, exactly as the substitution $\bar y:=\bar B$ does in~$\tau^\circ$.
\item Each filler $A'_s=\rho(A_s)$ contains no $p_{i_t}$, since $\rho$ renamed them all. It contains $r_t$ exactly where $A_s$ contains $p_{i_t}$, because $A_s$ contains no~$r_t$. All its other atoms are fixed by both $\rho$ and~$\sigma$. Hence $\sigma(A'_s)=A_s$.
\end{itemize}
So $\sigma(\iota')=\tau^\circ[\bar x:=\bar A,\ \bar y:=\bar B]=\iota$.

Finally, uniform substitution preserves classical validity. Suppose $\Pi'\models\varphi'$ and let~$v$ be any valuation. Define $v'(p):=v(\sigma(p))$. By induction on formulas, $v'(\psi)=v(\sigma\psi)$ for every~$\psi$. So if $v$ satisfies $\sigma\Pi'$, then $v'$ satisfies~$\Pi'$, hence~$\varphi'$, and $v$ satisfies $\sigma\varphi'$. Applied to $\iota'$, this shows that $\iota=\sigma(\iota')$ is sound.
\end{proof}

\begin{proof}[Proof of \cref{cor:search:purify}]
$\Rstar_{\CPC}$ is sound, so $A$ is sound iff every~$\tau_l$ is sound. By \cref{lem:search:purify}, this holds iff every~$\tau_l^\circ$ is sound, i.e.\ iff $A^\circ$ is sound.

Next, $A_1^\circ:=\bigcup_l\inst(\tau_l^\circ)$ is substitution-closed. Let $\tau^\circ\theta$ be an instance, where $\theta$ maps metavariables to formulas, and let $\sigma$ be a uniform substitution of formulas for atoms. Since $\tau^\circ$ contains no atom, $\sigma(\tau^\circ\theta)=\tau^\circ(\sigma\circ\theta)$, which is again an instance.

If $A^\circ$ is unsound, \cref{prop:search:post} gives $\Cl_{A^\circ}(\emptyset)=\Fm\ni\bot$. If $A^\circ$ is sound, then every step of~$A^\circ$ is in $\Sound(\Rstar_{\CPC})$ by completeness: $\Pi\models\varphi$ implies $\varphi\in\Cn_{\CPC}(\Pi)=\Cl_{\Rstar_{\CPC}}(\Pi)$. By \cref{lem:search:stepwise}, $\Cl_{A^\circ}(\emptyset)\subseteq\Cl_{\Rstar_{\CPC}}(\emptyset)=\Cn_{\CPC}(\emptyset)$, which does not contain~$\bot$.
\end{proof}

% ---------------------------------------------------------------------
\subsection{The commitment-order lemma}\label{app:search:commitment}

Recall \cref{def:search:scene}. An \emph{$R$-derivation} is a finite tree of sequents in which each node is either a leaf marked as a \emph{premise} (chosen freely by the prover) or, together with its children, forms an instance of some rule of~$R$: the children are the instance's premise sequents (none, for a premise-free rule) and the node is its conclusion sequent. The \emph{conclusion} of the derivation is the sequent at the root.

\begin{proof}[Proof of \cref{lem:search:commitment}(i)]
Let $G:=\{x:\rho_i\text{ is valid in }x\text{ for every }i\le m\}$. By the union bound over the whole library,
\[
\Pr_{x\sim D}[x\notin G]\ \le\ \sum_{i=1}^m\Pr_{x\sim D}[\rho_i\text{ is not valid in }x]\ \le\ \sum_{i=1}^m\varepsilon_i .
\]
Fix $x\in G$ and any $R$-derivation whose premises are all true in~$x$. We show by induction on height that every sequent in the tree is true in~$x$. Premise nodes are true by assumption. Every other node is the conclusion of an instance~$\iota$ of some~$\rho_i$, whose premise sequents are its children (possibly none) and are true in~$x$ by the induction hypothesis. Since $x\in G$, $\rho_i$ is valid in~$x$, so $\iota$ holds in~$x$ and its conclusion is true in~$x$. In particular the conclusion of the derivation is true in~$x$.

So the event ``some $R$-derivation from premises true in~$x$ has a conclusion false in~$x$'' is contained in the complement of~$G$. Its probability (its outer probability, if it is not measurable) is therefore at most $\sum_i\varepsilon_i$. The event~$G$ is determined by~$x$ and the library~$R$ alone. Because $R$ is fixed before $x$ is drawn, $G$ does not depend on the prover, and the bound holds simultaneously for every prover, including one that sees~$x$ and chooses any number of derivations of any length. The sum runs over the whole library, not over the rules a particular derivation uses.
\end{proof}

\begin{remark}[Per-pair guarantees do not give per-scene validity]\label{rem:app:search:perpair}
Suppose each scene in the support of~$D$ falsifies exactly one of $M$ instances of a rule~$\rho$, and the instance is drawn uniformly from these~$M$. Then $\Pr_{(x,\iota)}[\iota\text{ holds in }x]=1-1/M$, but $\rho$ is valid in no scene, so its per-scene validity is~$0$. A prover that sees~$x$ submits the falsified instance; this is regime~(ii) inside each scene.
\end{remark}

\begin{proof}[Proof of \cref{lem:search:commitment}(ii)]
Let $\iota_0$ be an instance of~$\rho$ that fails in~$x_0$: its premise sequents are true in~$x_0$ and its conclusion sequent is false in~$x_0$. The tree consisting of~$\iota_0$ alone (its conclusion at the root, its premises as leaves) is a $\{\rho\}$-derivation from premises true in~$x_0$ with a conclusion false in~$x_0$. A prover that knows~$x_0$ submits it, and nothing in this regime is random, so this happens with probability~$1$. If $\mu\{\iota:\iota\text{ fails in }x_0\}>0$, such an~$\iota_0$ exists. If $\mu$ is supported on instances that hold in~$x_0$, the average validity is~$1$, i.e.\ $\varepsilon=0$, but a failing instance off the support of~$\mu$ may still exist, and the same prover uses it. So no bound in terms of~$\varepsilon$ holds.
\end{proof}

For part~(iii) we use the Schwartz--Zippel lemma \citep{schwartz1980fast,zippel1979probabilistic,demillo1978probabilistic}: if $f\in F[z_1,\dots,z_n]$ is nonzero of total degree at most~$d$ and $S\subseteq F$ is finite, then $\Pr_{a\sim\mathrm{Unif}(S^n)}[f(a)=0]\le d/|S|$.

\begin{proof}[Proof of \cref{lem:search:commitment}(iii)]
Formalize the protocol as follows. In round~$k$ the prover, which may be randomized, adaptive and computationally unbounded, commits to explicit expressions for $p_k,q_k\in F[z_1,\dots,z_n]$. The verifier computes an upper bound on $\deg(p_k-q_k)$ from the expressions and rejects if it exceeds~$d$. Otherwise it draws $a_k$ uniformly from $S_k^n$, independently of everything so far, and accepts iff $p_k(a_k)=q_k(a_k)$.

Let $\mathcal H_k$ be the $\sigma$-algebra generated by everything that happens before $a_k$ is drawn: earlier claims, points and answers, and the prover's coins. The claim $(p_k,q_k)$ is $\mathcal H_k$-measurable, and $a_k$ is independent of~$\mathcal H_k$. Let $E_k$ be the event that claim~$k$ passes the degree check and $p_k-q_k$ is nonzero in $F[z_1,\dots,z_n]$, and $\mathrm{Acc}_k$ the event that it is accepted. By the Schwartz--Zippel lemma, applied conditionally on~$\mathcal H_k$,
\[
\Pr[E_k\cap\mathrm{Acc}_k\mid\mathcal H_k]\ =\ \mathbf 1_{E_k}\cdot\Pr_{a\sim\mathrm{Unif}(S_k^n)}\big[(p_k-q_k)(a)=0\big]\ \le\ d/|S_k| .
\]
Taking expectations, $\Pr[E_k\cap\mathrm{Acc}_k]\le d/|S_k|$. A false claim (one with $p_k-q_k\neq0$ over~$F$) that fails the degree check is rejected, so a false claim is accepted only on $E_k\cap\mathrm{Acc}_k$. By the union bound,
\[
\Pr\big[\exists k\le N:\ \text{claim $k$ is false and accepted}\big]\ \le\ \sum_{k\le N}d/|S_k| ,
\]
which is $Nd/|S|$ when $S_k=S$ for all~$k$. If $|S_k|\ge d\,2^k/\delta$ for all $k\ge1$, then $\sum_{k\ge1}d/|S_k|\le\sum_{k\ge1}\delta2^{-k}=\delta$, which bounds the probability that any false claim in an unbounded sequence is accepted. The prover's strategy enters only through the~$\mathcal H_k$, so the bounds hold uniformly over provers.
\end{proof}

\paragraph{Why the conditions of (iii) are needed.}
\begin{itemize}
\item \emph{Degree enforcement.} Over $F=\mathbb F_{2^k}$ with $S=F$, the claim $z^{2^k}\equiv z$ can be written with $k$ squarings. The polynomial $z^{2^k}-z$ is nonzero of degree~$2^k$, but it vanishes on all of~$F$, so random evaluation always accepts. The bound $d/|S|$ is vacuous here, since $d=2^k=|S|$. A verifier that read the degree off the length of the expression, roughly~$k$, would wrongly believe its error to be about $k/2^k$.
\item \emph{The field.} For prime~$p$, $f=(a+b)^p-a^p-b^p=\sum_{0<i<p}\binom pi a^ib^{p-i}$ is nonzero over~$\mathbb Z$, but $p$ divides every $\binom pi$ with $0<i<p$, so $f=0$ in $\mathbb F_p[a,b]$. Evaluation mod~$p$ accepts this false integer identity with probability~$1$. Over~$\mathbb Q$ with $S\subseteq\mathbb Z$, the lemma applies with $d=p$.
\item \emph{Order of commitment.} If the evaluation points were drawn before the prover committed to its claims, and visible to it, the prover could choose a nonzero $p-q$ that vanishes at those points. This is regime~(ii) again.
\end{itemize}

% ---------------------------------------------------------------------
\subsection{The doctrinal paradox}\label{app:search:doctrinal}

\begin{proof}[Proof of \cref{thm:search:doctrinal}]
\emph{Step 1: each $h_i$ is a coherent, cut-closed calculus.} The sets $T_1,T_2,T_3$ are satisfied by the valuations $(v(p),v(q))=(1,1),(1,0),(0,1)$ respectively. For every set~$\Gamma$, $\Cl_{h_i}(\Gamma)=\Cn_{\CPC}(T_i\cup\Gamma)$. For~$\supseteq$: if $T_i\cup\Gamma\models\varphi$, compactness gives a finite $\Gamma_0\subseteq\Gamma$ with $T_i\cup\Gamma_0\models\varphi$, so $\Gamma_0\step\varphi\in h_i$ and $\varphi\in\Cl_{h_i}(\Gamma)$. For~$\subseteq$: $\Cn_{\CPC}(T_i\cup\Gamma)$ contains~$\Gamma$ and is closed under~$h_i$, since $\Pi\subseteq\Cn_{\CPC}(T_i\cup\Gamma)$ and $T_i\cup\Pi\models\varphi$ imply $T_i\cup\Gamma\models\varphi$. In particular $\Cl_{h_i}(\emptyset)=\Cn_{\CPC}(T_i)$, which does not contain~$\bot$ because $T_i$ is satisfiable. So every~$h_i$ is $\Ac$-coherent and is a possible target.

\emph{Step 2: memberships.} The following table records which $h_i$ contain each step of~$\pi$.
\begin{center}\small
\begin{tabular}{@{}lccc@{}}
\toprule
step & $h_1$ ($p,q$) & $h_2$ ($p,\neg q$) & $h_3$ ($\neg p,q$)\\
\midrule
$\step p$ & $\checkmark$ & $\checkmark$ & -- \\
$\step q$ & $\checkmark$ & -- & $\checkmark$ \\
$p,q\step p\wedge q$ & $\checkmark$ & $\checkmark$ & $\checkmark$ \\
$\step\neg(p\wedge q)$ & -- & $\checkmark$ & $\checkmark$ \\
$p\wedge q,\ \neg(p\wedge q)\step\bot$ & $\checkmark$ & $\checkmark$ & $\checkmark$ \\
\bottomrule
\end{tabular}
\end{center}
The checkmarks are immediate: $p\in T_1,T_2$; $q\in T_1,T_3$; $\neg q\models\neg(p\wedge q)$ and $\neg p\models\neg(p\wedge q)$; and the third and fifth steps are valid outright. For the dashes: $T_3\models\neg p$, $T_2\models\neg q$ and $T_1\models p\wedge q$, and each $T_i$ is satisfiable, so $T_3\not\models p$, $T_2\not\models q$ and $T_1\not\models\neg(p\wedge q)$.

\emph{Step 3: $\pi$ is a legal detection that deletes nothing.} The five steps form an argument~$\pi$ from~$\emptyset$ to~$\bot$: the first two yield $p$ and~$q$, the third yields $p\wedge q$, the fourth yields $\neg(p\wedge q)$, and the fifth yields~$\bot$. With uniform weights, $s$ is in the majority set over~$\Hc$ iff $s$ lies in at least two of the~$h_i$, since $2/3>1/2>1/3$. By the table, every step of~$\pi$ lies in at least two $h_i$, so $\mathrm{Steps}(\pi)\subseteq\hat R_{\mathrm{maj}}$, the majority set over~$\Hc$. Also by the table, $\mathrm{Steps}(\pi)\not\subseteq h_i$ for each~$i$. So the deletion set $\{R\in\Hc:\mathrm{Steps}(\pi)\subseteq R\}$ is empty.

\emph{Step 4: induction on rounds.} Fix any target $h_{i^*}$. Consider the environment that, in each round, exhibits~$\pi$ or presents the positive datum $\step p\vee\neg p$, in any order. We show $\VS_t=\Hc$ for all~$t$. Initially $\VS_1=\Hc$. If $\VS_t=\Hc$, then $\hat R_t=\hat R_{\mathrm{maj}}\supseteq\mathrm{Steps}(\pi)$. So $\pi$ is a legal (N)-datum, since $\emptyset\in\Ac$, and by Step~3 it deletes nothing. The datum $\step p\vee\neg p$ lies in every~$h_i$, so it is a legal (P)-datum for the target and deletes $\{R:\step p\vee\neg p\notin R\}=\emptyset$. Hence $\VS_{t+1}=\Hc$. So $\VS_t=\Hc$ and $\hat R_t=\hat R_{\mathrm{maj}}$ for all~$t$, and $\pi$ is available in every round.
\end{proof}

\paragraph{Consequences recorded in the main text.}
\begin{itemize}
\item \emph{The vote is unsound relative to every target.} By Step~3, $\bot\in\Cl_{\hat R_{\mathrm{maj}}}(\emptyset)$. By Step~1, $\bot\notin\Cl_{h_i}(\emptyset)$ for each~$i$. By \cref{lem:search:stepwise} with $B=\emptyset$, $\hat R_{\mathrm{maj}}\not\subseteq\Sound(h_i)$ for each~$i$.
\item \emph{Discriminating data break the paradox.} Let the target be~$h_1$ and present $\step p\in h_1$. This deletes~$h_3$, since $\step p\notin h_3$. Over $\VS=\{h_1,h_2\}$, with weight $1/3$ each, a step is a strict majority iff it lies in both, so $\hat R=h_1\cap h_2$. Since $\step q\notin h_2$, $\pi$ is no longer available. The $h_i$ are pairwise incomparable: $\step p\wedge q$ lies only in~$h_1$, $\step p\wedge\neg q$ only in~$h_2$, and $\step\neg p\wedge q$ only in~$h_3$. So on a text enumerating all of the target, each wrong hypothesis is eventually deleted by a datum it lacks, and from then on the uninformative detections stop.
\end{itemize}
